\documentclass[11pt]{article}

\usepackage[a4paper,margin=1in]{geometry}
\usepackage{amsmath,amssymb,amsthm,mathtools}
\usepackage{array}
\usepackage{booktabs}
\usepackage{enumitem}
\usepackage{tikz}
\usepackage{physics}
\usepackage{needspace}
\usepackage{xcolor}
\definecolor{courseblue}{RGB}{0,90,150}
\usetikzlibrary{arrows.meta,positioning}

\setlength{\parindent}{0pt}
\setlength{\parskip}{0.75em}

\newcounter{problem}
\newenvironment{problem}{%
    \refstepcounter{problem}%
    \Needspace{7\baselineskip}%
    \par\vspace{0.7em}%
    \noindent\textcolor{courseblue}{\rule{\textwidth}{0.4pt}}%
    \par\vspace{0.5em}%
    \begin{list}{}{%
        \setlength{\leftmargin}{2.3em}%
        \setlength{\labelwidth}{1.8em}%
        \setlength{\labelsep}{0.5em}%
        \setlength{\topsep}{0pt}%
        \setlength{\partopsep}{0pt}%
    }%
    \item[{[\theproblem]}]
}{%
    \end{list}
    \par\vspace{0.5em}
}

\newenvironment{parts}{%
    \begin{enumerate}[label={[\alph*]},leftmargin=2.2em,itemsep=0.65em,topsep=0.65em]
}{%
    \end{enumerate}
}

% Type each solution between \begin{answer} and \end{answer}.
\newenvironment{answer}{%
    \Needspace{4\baselineskip}%
    \par\vspace{0.35em}
    \noindent\textbf{Answer:}\par
}{%
    \par\vspace{0.25em}
}



\title{CO250 Assignment 3}
\author{}
\date{September 28--October 5}

\begin{document}
\maketitle

\begin{problem}
    For each of the following linear programs determine whether it is \textbf{infeasible}, \textbf{unbounded}, or
    if it has an \textbf{optimal point}. In each case, justify your answer (e.g.\ with a certificate that the
    linear program has this property).

    \begin{parts}
        \Needspace{12\baselineskip}\item
        \[
        \begin{aligned}
            \min_{x_1,x_2}\quad & x_1+2x_2 \\
            \text{s.t.}\quad & -x_1\leq 0, \\
            & -x_2\leq 0, \\
            & -x_1-x_2\leq -1.
        \end{aligned}
        \]
        \begin{answer}
            Unbounded. Take $\bar{x} = (1,1)$ and $d = (1,1)$.
        \end{answer}

        \Needspace{13\baselineskip}\item Consider
        \[
        \begin{aligned}
            \min_{x_1,x_2}\quad & -x_1-x_2 \\
            \text{s.t.}\quad & x_1\leq 2, \\
            & x_2\leq 3, \\
            & -x_1\leq 0, \\
            & -x_2\leq 0.
        \end{aligned}
        \]
        \begin{answer}

        \end{answer}

        \Needspace{11\baselineskip}\item
        \[
        \begin{aligned}
            \min_{x_1,x_2}\quad & -x_1+x_2 \\
            \text{s.t.}\quad & -x_1\leq 0, \\
            & -x_2\leq 0.
        \end{aligned}
        \]
        \begin{answer}

        \end{answer}

        \Needspace{11\baselineskip}\item
        \[
        \begin{aligned}
            \min_{x_1,x_2}\quad & x_1+x_2 \\
            \text{s.t.}\quad & x_1\leq 0, \\
            & -x_1\leq -1.
        \end{aligned}
        \]
        \begin{answer}

        \end{answer}
    \end{parts}
\end{problem}

\begin{problem}
    A company has a set $J$ of jobs to be assigned to a set $I$ of employees. Employee $i\in I$ needs $c_{ij}$
    hours to complete job $j\in J$. The company wants to complete all jobs by assigning exactly one
    job to each employee, while minimizing the total number of hours spent to complete all jobs.
    For this reason, the company has formulated the following IP:
    \[
    \begin{gathered}
        \min \sum_{i\in I}\sum_{j\in J}c_{ij}x_{ij} \\[0.5em]
        \begin{aligned}
            \sum_{i\in I}x_{ij} &= 1 && \forall j\in J \\
            \sum_{j\in J}x_{ij} &= 1 && \forall i\in I
        \end{aligned} \\[0.5em]
        x\geq\mathbf{0},\text{ integer}
    \end{gathered}
    \]

    Consider the LP obtained by the above IP after ignoring the integrality constraints, that is:
    \[
    \begin{gathered}
        \min \sum_{i\in I}\sum_{j\in J}c_{ij}x_{ij} \\[0.5em]
        \begin{aligned}
            \sum_{i\in I}x_{ij} &= 1 && \forall j\in J \\
            \sum_{j\in J}x_{ij} &= 1 && \forall i\in I
        \end{aligned} \\[0.5em]
        x\geq\mathbf{0},
    \end{gathered}
    \]

    Show that the above LP is feasible (i.e., that it admits at least one feasible solution) if and
    only if $|I|=|J|$.

    \textbf{Note:} to show that the LP is feasible, you need to provide one feasible solution. To show that
    the LP is infeasible, provide a certificate of infeasibility as in Proposition 2.1 (p.\ 46) in the
    textbook.
    \begin{answer}

    \end{answer}
\end{problem}

\Needspace{22\baselineskip}
\begin{problem}
    \begin{parts}
        \Needspace{16\baselineskip}\item Here is a linear program
        \[
        \begin{aligned}
            \text{minimize}\quad & 3x_1+2x_2+x_3 \\
            \text{subject to}\quad & x_1+x_2+x_3\leq 8, \\
            & 2x_1-x_2+x_3\geq 2, \\
            & x_1+3x_2-x_3\leq 10, \\
            & -x_1+x_2+2x_3\geq 1.\quad x_1,x_2,2x_3\geq 0.
        \end{aligned}
        \]
        % The source prints x_1, x_2, 2x_3 >= 0; preserved verbatim.
        Write an equivalent linear program in standard equality form.
        \begin{answer}

        \end{answer}

        \Needspace{17\baselineskip}\item Consider the following matrix (where the top row just numbers the columns).
        \[
        A=\bordermatrix{
            & 1 & 2 & 3 & 4 & 5 \cr
            & 1 & 0 & 1 & 2 & 1 \cr
            & 0 & 1 & 1 & 1 & 2 \cr
            & 0 & 0 & 0 & 1 & 1 \cr
        }.
        \]
        which of the following sets of columns are bases of this matrix?
        \[
            \{1,2,5\},\ \{1,2\},\ \{1,2,3\},\ \{1,3,5\}
        \]
        \begin{answer}

        \end{answer}

        \Needspace{13\baselineskip}\item Consider the same matrix as in (b) in the equation
        \[
            Ax=\begin{pmatrix}1\\2\\3\end{pmatrix}.
        \]
        Put this linear equation in canonical form, using the basis $\{1,2,4\}$.
        \begin{answer}

        \end{answer}
    \end{parts}
\end{problem}

\end{document}
